\section{Conclusion}

To what extent can NLP models detect actionable improvement
suggestions within high-star hotel reviews? Our experiments show
that suggestion mining is feasible in the hotel review domain, but
cross-domain transfer from software forums degrades substantially
without domain adaptation. BERT stage-2, fine-tuned on just 320
hotel-domain sentences, achieves suggestion-class F1 of 0.581 on
the MBS test split---approximately double the cross-domain-only
result (0.286) and the best performance among all models tested.

Our contributions are:
\begin{enumerate}
  \item Domain-adapted annotation guidelines for hotel review
  suggestions, validated with Fleiss' $\kappa$ of 0.379 on a
  100-sentence calibration sample across four annotators.
  \item A three-tier model comparison demonstrating that
  contextual representations (BERT) outperform bag-of-words
  features (TF-IDF~+~LR, suggestion-class F1 0.596) and surface
  patterns (regex, 0.386) on in-domain data, but that this
  advantage does not automatically transfer cross-domain.
  \item Two-stage fine-tuning that improves MBS suggestion-class
  F1 by 0.295 points (from 0.286 to 0.581), recovering recall
  from 0.200 to 0.600 with only 320 in-domain training sentences.
  \item Thematic analysis of 4,179 predicted suggestions (from
  BERT stage-2, suggestion-class F1 = 0.581) indicating that
  room-related suggestions account for the largest share
  (23.7\%), followed by facilities (15.3\%) and food and
  beverage (11.7\%). At precision 0.563, roughly 44\% of inputs
  to BERTopic may be false positives, so these distributions are
  approximate.
\end{enumerate}

Future work should scale annotation to more hotels and review
platforms, investigate prompt-based suggestion detection using
large language models, and explore multi-label classification of
suggestion types (explicit, implicit, comparative).
